\documentclass[12pt,openany]{book}

\usepackage{amsmath,amsthm,amsfonts,amscd} % Packages for mathematics

% Colors
\usepackage[dvipsnames]{xcolor}
\definecolor{titleblue}{RGB}{0,53,128}
\definecolor{chaptergray}{RGB}{140,140,140}
\definecolor{sectiongray}{RGB}{180,180,180}

\definecolor{thmcolor}{RGB}{231, 76, 60}
\definecolor{defcolor}{RGB}{52, 152, 219}
\definecolor{lemcolor}{RGB}{155, 89, 182}
\definecolor{corcolor}{RGB}{46, 204, 113}
\definecolor{procolor}{RGB}{241, 196, 15}

% Fonts
\usepackage[T1]{fontenc}
\usepackage[utf8]{inputenc}
\usepackage{newpxtext,newpxmath}
\usepackage{sectsty}
\allsectionsfont{\sffamily\color{titleblue}\mdseries}

% Page layout
\usepackage{geometry}
\geometry{a4paper,left=1.325in,right=1in,top=1in,bottom=1in,heightrounded}
\usepackage{fancyhdr}
\fancyhf{}
\fancyhead[LE,RO]{\thepage}
\fancyhead[LO]{\nouppercase{\rightmark}}
\fancyhead[RE]{\nouppercase{\leftmark}}
\renewcommand{\headrulewidth}{0.5pt}
\renewcommand{\footrulewidth}{0pt}

% Chapter formatting
\usepackage{titlesec}
\titleformat{\chapter}[display]
{\normalfont\sffamily\Huge\bfseries\color{titleblue}}{\chaptertitlename\ \thechapter}{20pt}{\Huge}
\titleformat{\section}
{\normalfont\sffamily\Large\bfseries\color{titleblue!100!gray}}{\thesection}{1em}{}
\titleformat{\subsection}
{\normalfont\sffamily\large\bfseries\color{titleblue!75!gray}}{\thesubsection}{1em}{}

% Table of contents formatting
\usepackage{tocloft}
\renewcommand{\cftchapfont}{\sffamily\color{titleblue}\bfseries}
\renewcommand{\cftsecfont}{\sffamily\color{chaptergray}}
\renewcommand{\cftsubsecfont}{\sffamily\color{sectiongray}}
\renewcommand{\cftchapleader}{\cftdotfill{\cftdotsep}}

% Hyperlinks
\usepackage{hyperref}
\hypersetup{
	colorlinks=true,
	linkcolor=titleblue,
	filecolor=black,      
	urlcolor=titleblue,
}

%Listing
\usepackage{listings} %Code
\renewcommand{\lstlistingname}{Code}%

\definecolor{sagegreen}{rgb}{0.0,0.6,0.4}
\definecolor{sagepurple}{rgb}{0.6,0.0,0.4}
\definecolor{sageblue}{rgb}{0.0,0.4,0.6}
\definecolor{sageorange}{rgb}{1.0,0.4,0.0}
\definecolor{sagegray}{rgb}{0.4,0.4,0.4}

\lstdefinestyle{sage}{
	language=Python,
	backgroundcolor=\color{white},
	basicstyle=\small\ttfamily\color{black}, 
	basicstyle=\footnotesize\ttfamily\color{black},
	keywordstyle=\color{blue!60!black},
	commentstyle=\color{green!60!black},
	stringstyle=\color{purple!60!black},
	showstringspaces=false,
	breaklines=true,
	tabsize=4,
	morekeywords={True, False, None},
	frame=leftline, % Remove the border
	framesep=3pt,
	frameround=tttt,
	framexleftmargin=3pt,
	numbers=left,
	numberstyle=\small\color{gray},
	xleftmargin=15pt, % Increase the left margin
	xrightmargin=5pt,
	captionpos=b,
	belowskip=0pt,
	aboveskip=4pt
}

%Ceiling and Floor Function
\usepackage{mathtools}
\DeclarePairedDelimiter{\ceil}{\lceil}{\rceil}
\DeclarePairedDelimiter{\floor}{\lfloor}{\rfloor}

%Algorithm
\usepackage[ruled,linesnumbered]{algorithm2e}
\usepackage{setspace}
\usepackage{algpseudocode}
\SetKwComment{Comment}{/* }{ */}
\SetKw{Break}{break}
\SetKw{Downto}{downto}
\SetKwProg{Fn}{Function}{:}{end}
\SetKwFunction{KeyGen}{KeyGen}


%---------------------------My Preamble
\usepackage{marvosym} %Lightning
\usepackage{booktabs}
\usepackage{multicol}
\setlength{\columnsep}{2cm}
\setlength{\columnseprule}{1.25pt}
\usepackage{enumerate}
\usepackage{soul}
\newcommand{\mathcolorbox}[2]{\colorbox{#1}{$\displaystyle #2$}}
\usepackage{graphicx}
\usepackage{tikz}
\usepackage{tikz-cd}
\usetikzlibrary{calc}
\usetikzlibrary{arrows, arrows.meta, positioning, shapes.multipart}

%Tcolorbox
\usepackage[most]{tcolorbox}
\tcbset{colback=white, arc=5pt}
%\tcbset{enhanced, colback=white,colframe=black,fonttitle=\bfseries,arc=4mm,boxrule=1pt,shadow={2mm}{-1mm}{0mm}{black!50}}
%White box with black text and shadow
%\begin{tcolorbox}[colback=white,colframe=black,fonttitle=\bfseries,title=Black Shadow Box,arc=4mm,boxrule=1pt,shadow={2mm}{-1mm}{0mm}{black!50}]
%	This is a white box with black text and a subtle shadow. The shadow adds some depth and dimension to the box without overpowering the design.
%\end{tcolorbox}

%Theorem
\newtheorem{axiom}{Axiom}[chapter]
\newtheorem{theorem}{Theorem}[chapter]
\newtheorem{proposition}[theorem]{Proposition}
\newtheorem{corollary}{Corollary}[theorem]
\newtheorem{lemma}[theorem]{Lemma}

\theoremstyle{definition}
\newtheorem{definition}{Definition}[chapter]
\newtheorem{remark}{Remark}[chapter]
\newtheorem{exercise}{Exercise}[chapter]
\newtheorem{example}{Example}[chapter]
\newtheorem*{note}{Note}

%New Command
\newcommand{\set}[1]{\left\{#1\right\}}
\newcommand{\N}{\mathbb{N}}
\newcommand{\Z}{\mathbb{Z}}
\newcommand{\Q}{\mathbb{Q}}
\newcommand{\R}{\mathbb{R}}
\newcommand{\C}{\mathbb{C}}
\newcommand{\F}{\mathbb{F}}

% Begin document
\begin{document}
	
	% Title page
	\begin{titlepage}
		\begin{center}
			{\Huge\textsf{\textbf{Advanced Applied Mathematics}}\par}
			\vspace{0.5in}
			{\Large Ji Yong-Hyeon\par}
			\vspace{1in}
			\includegraphics[scale=.5]{aam.png}\par
			\vspace{1in}
			{Department of Information Security, Cryptology, and Mathematics\par}
			{College of Science and Technology\par}
			{Kookmin University\par}
			%\includegraphics[width=1.5in]{school_logo.jpg}\par
			\vspace{.25in}
			{\large \today\par}
		\end{center}
	\end{titlepage}
	
	\iffalse
	% Title page
	\begin{titlepage}
		\begin{center}
			{\Huge\textsf{\textbf{Abstract Algebra}}\par}
			\vspace{0.5in}
			{\Large Ji Yong-Hyeon\par}
			\vspace{1in}
			\includegraphics[width=3in]{algebra.png}\par
			\vspace{1in}
			{\large Ji Yong-Hyeon\par}
			{\large \today\par}
		\end{center}
	\end{titlepage}
	\fi
	
	% Table of contents
	\tableofcontents
	=m
	% Chapters
	\mainmatter
	
	\chapter{Linear Algebra}
	\section{System of Linear Equations}
	\begin{itemize}
		\item A system of linear equations
		\[
		\begin{cases}
			x_1,\dots,x_n:\text{unknowns}\\
			\text{\# of unknowns}=n\\
			\text{\# of equations}=m
		\end{cases}
		\]
		\begin{align*}
			&\begin{cases}
				a_{11}x_1+a_{12}x_2+\cdots+a_{1n}x_n=b_1\\
				a_{21}x_1+a_{22}x_2+\cdots+a_{2n}x_n=b_2\\
				\hspace{50pt}\vdots\\
				a_{m1}x_1+a_{m2}x_2+\cdots+a_{mn}x_n=b_m\\
			\end{cases}\\
			\iff&\begin{bmatrix}
				a_{11} & a_{12} & \cdots & a_{1n}\\
				\vdots \\
				a_{m1} & a_{m2} &\cdots & a_{mn}
			\end{bmatrix}\begin{bmatrix}
				x_1\\ \vdots \\ x_n
			\end{bmatrix}=\begin{bmatrix}
				b_1\\ \vdots \\ b_m
			\end{bmatrix} &\textcolor{blue}{A\textbf{x}=\textbf{b}}\\
			\iff& x_1\begin{bmatrix}
				a_{11}\\ a_{21}\\ \vdots\\ a_{m1}
			\end{bmatrix}+x_2\begin{bmatrix}
				a_{12}\\ a_{22}\\ \vdots\\ a_{m2}
			\end{bmatrix}+\cdots+x_n\begin{bmatrix}
				a_{1m}\\ a_{2m}\\ \vdots\\ a_{mn}
			\end{bmatrix}=\begin{bmatrix}
				b_1\\ b_2\\ \vdots\\ b_{m} 
			\end{bmatrix} & \textcolor{blue}{x_1\textbf{C}_1+\cdots+x_n\textbf{C}_n=\textbf{b}}
		\end{align*}
		\item Matrix operation
		\begin{enumerate}[(i)]
			\item scalar multiplication: $kA$
			\item addition: $A+B$
			\item multiplication: $AB$
		\end{enumerate}
		\item Properties \begin{itemize}
			\item Associative: $(A+B)+C=A+(B+C)$, $A(BC)=(AB)C$
			\item Distributive: $(AB)C=A(BC)$
			\item (in general) not commutative: $AB\neq BA$
		\end{itemize}
		\item Transpose of $A$: $A^T$ \[
		(a_{ij})_{m\times n}\longrightarrow(a^t_{ij})_{n\times m}=(a_{ji})_{n\times m}
		\]
		\item Square Matrices
	\end{itemize}
	
	\subsection{Elementary Transformation}
	\begin{itemize}
		\item Exchange of two equations (rows in the matrix representing the system
		of equations)
		\item Multiplication of an equation (row) with a constant $\lambda\in\R^*$
		\item Addition of two equations (rows)
	\end{itemize}
	
	\begin{remark}
		$A\textbf{x}=\textbf{b}\iff[A\mid\textbf{b}]$.
	\end{remark}
	
	\begin{example}
		\begin{align*}
			\begin{bmatrix}
				1 & 1 & 1\\ 1 & -1& 2\\ 2 & 0 & 3
			\end{bmatrix}\begin{bmatrix}
				x_1\\ x_2\\ x_3
			\end{bmatrix}=\begin{bmatrix}
				3 \\ 2 \\ 5
			\end{bmatrix} &\iff \left[
			\begin{array}{ccc|c}
				1 & 1 & 1 & 3\\
				1 & -1 & 2 & 2\\
				2 & 0 & 3 & 5
			\end{array}
			\right] \\ &\xLeftrightarrow[R_3\gets R_3-2R_1]{R_2\gets R_2-R_1} \left[
			\begin{array}{ccc|c}
				1 & 1 & 1 & 3\\
				0 & -2 & 1 & -1\\
				0 & -2 & 1 & -1
			\end{array} 
			\right] \\
			&\xLeftrightarrow[R_2\gets -\frac{1}{2}R_2]{R_3\gets R_3-R_2} \left[
			\begin{array}{ccc|c}
				1 & 1 & 1 & 3\\
				0 & 1 & -1/2 & 1/2\\
				0 & 0 & 0 & 0
			\end{array} 
			\right]\quad\text{Row-Echelon Form (REF)} \\
			&\xLeftrightarrow{R_1\gets R_1-R_2} \left[
			\begin{array}{ccc|c}
				1 & 0 & 3/2 & 5/2\\
				0 & 1 & -1/2 & 1/2\\
				0 & 0 & 0 & 0
			\end{array} 
			\right]\quad\text{Reduced Row-Echelon Form (RREF)} \\
			&\iff \begin{cases}
				x_1=-\frac{3}{2}x_3+\frac{5}{2}\\
				x_2=\frac{1}{2}x_3+\frac{1}{2}
			\end{cases}.
		\end{align*} Let $x_3=\lambda$ then \[
		\textbf{x}=\begin{bmatrix}
			-\frac{3}{2}\lambda+\frac{5}{2}\\ \frac{1}{2}\lambda+\frac{1}{2}\\ \lambda
		\end{bmatrix}=\begin{bmatrix}
		\frac{5}{2}\\ \frac{1}{2}\\ 0
		\end{bmatrix}+\lambda\begin{bmatrix}
			-\frac{3}{2}\\ \frac{1}{2}\\ 1
		\end{bmatrix}.
		\]
	\end{example}

	\section{Reduced Row-Echelon Form (RREF)}
	\section{General Solution and Minus-1 Trick}
	
	\section{Vector Space, Linear Transformation, Basis}
	
	\chapter{Analytic Geometry}
	\section{Rank-Nullity Theorem}
	\section{Positive Definite, Gram-Schmidt Process}
	
	\chapter{Matrix Decompositions}
	\section{Eigenvalues and Eigenvectors}
	\section{Complex Matrix, Hermitian}
	\section{Decomposition of Matrices: Spectral Decomposition, SVD}
	
	\chapter{Vector Calculus}
	\section{Vector Calculus}
	
	\chapter{Probability and Distributions}
	\section{Backpropagation, Probability Theory Basic}
	\section{Gaussian Distribution and its Applications}
	\section{Continuous Optimization}
	\section{Probabilistic Modeling and Inference}
	\section{Linear Regression}
	\section{Supervised Learning}
	\section{Bayesian Linear Regression}
	\section{Maximum Likelihood Estimation(MLE)}
	\section{Generalized Linear Regression}
	\section{Maximum A Posteriori Estimation(MAP)}
	\section{Principal Component Analysis(PCA)}
	
	\chapter{Continuous Optimization}
	
	\chapter{When Models Meet Data}
	
	\chapter{Linear Regression}
	
	\chapter{Dimensionality Reduction with PCA}
	% End document
\end{document}
